% Figure: the dynamic vibration absorber. The primary-mass response
% magnitude versus forcing frequency, comparing the bare primary system
% to the same system fitted with a tuned absorber. The single resonant
% peak of the bare system splits into two peaks straddling an antiresonance
% where the absorber holds the primary mass still.
\begin{figure}[htbp]
\centering
\begin{tikzpicture}
\begin{axis}[
  width=12cm, height=7cm,
  xlabel={$\omega/\omega_n$ (primary, undamped)},
  ylabel={$|X_1|\,/\,(F_0/k_1)$},
  xmin=0.6, xmax=1.5, ymin=0, ymax=10.5,
  axis lines=left,
  domain=0.6:1.5, samples=600,
  legend pos=north east,
  legend cell align=left,
  every axis x label/.style={at={(current axis.right of origin)}, anchor=west},
  every axis y label/.style={at={(current axis.above origin)}, anchor=south},
  grid=major, grid style={enggray!20},
  tick label style={font=\scriptsize},
  label style={font=\small},
  legend style={font=\scriptsize},
  restrict y to domain=0:60,
]
% Bare primary system, lightly damped: 1/sqrt((1-r^2)^2+(2 zeta1 r)^2)
\addplot[enggray, thick, dashed] {1/sqrt((1-x^2)^2 + (2*0.02*x)^2)};
\addlegendentry{bare primary ($\zeta_1=0.02$)}
% With tuned absorber: mass ratio mu=0.05, tuned f=1, absorber damping zeta_a.
% Magnitude of primary response for 2-DOF absorber, mu=0.05, beta=1.
% |X1 k1 / F0| = | (beta^2 - r^2) | / |det|, here we plot a representative
% two-peak curve from the standard absorber transfer function with small
% absorber damping za=0.06. Use the closed form with mu, beta, za.
% Numerator: sqrt( (beta^2 - r^2)^2 + (2 za beta r)^2 )
% Denominator: sqrt( [ (1 - r^2)(beta^2 - r^2) - mu beta^2 r^2 ]^2
%                   + (2 za beta r)^2 [ 1 - r^2 (1+mu) ]^2 )
\addplot[engblue, very thick] {
  sqrt( (1 - x^2)^2 + (2*0.06*1*x)^2 )
  /
  sqrt(
    ( (1 - x^2)*(1 - x^2) - 0.05*1*x^2 )^2
    + (2*0.06*1*x)^2 * ( 1 - x^2*(1+0.05) )^2
  )
};
\addlegendentry{with absorber ($\mu=0.05$, $\zeta_a=0.06$)}
\node[engred, font=\scriptsize, anchor=south] at (axis cs:1.0,0.7) {antiresonance};
\draw[engred, thin, ->] (axis cs:1.0,0.6) -- (axis cs:1.0,0.15);
\end{axis}
\end{tikzpicture}
\caption[The dynamic vibration absorber]{Steady-state amplitude of the
primary mass, normalised to its static deflection, against forcing
frequency. The bare lightly-damped primary system (grey dashed) has a
single tall resonance near $\omega/\omega_n=1$. Adding a tuned absorber
(an auxiliary mass on its own spring, tuned to the same frequency, mass
ratio $\mu=0.05$) splits the single peak into two lower peaks straddling
an antiresonance at the tuning frequency, where the absorber's reaction
force cancels the applied force and the primary mass is held nearly
still. A small amount of absorber damping ($\zeta_a=0.06$) caps the two
new peaks; the undamped absorber would drive the primary response to
zero exactly at tuning but leave two infinite peaks on either side. The
absorber trades one large resonance for two smaller ones and a quiet
band in between, the principle behind the tuned mass damper of a tall
building and the harmonic balancer of an engine.}
\label{fig:vol03ch06:absorber}
\end{figure}
